\documentclass[11pt]{article}

%% Page geometry
\usepackage[a4paper, margin=1in]{geometry}

%% Core math + fonts
\usepackage{amsthm,amsmath,amsfonts,amssymb}
\usepackage{microtype}

%% Citations
\usepackage[authoryear, round]{natbib}

%% Hyperlinks
\usepackage[colorlinks, citecolor=blue, urlcolor=blue, linkcolor=blue]{hyperref}
\usepackage{xurl}

%% Graphics & figures
\usepackage{graphicx}
\graphicspath{{./images/}}
\usepackage{float}
\usepackage{subcaption}
\usepackage{booktabs}

%% Algorithms
\usepackage[linesnumbered, ruled, vlined]{algorithm2e}

%% TikZ for diagrams
\usepackage{tikz}
\usetikzlibrary{shapes,arrows,fit,calc,positioning}
\tikzset{box/.style={draw, diamond, thick, text centered, minimum height=0.5cm, minimum width=1cm}}
\tikzset{line/.style={draw, thick, -latex'}}

%% Misc
\usepackage{comment}
\usepackage{dsfont}

%% Theorem environments
\theoremstyle{plain}
\newtheorem{axiom}{Axiom}
\newtheorem{claim}[axiom]{Claim}
\newtheorem{theorem}{Theorem}[section]
\newtheorem{lemma}[theorem]{Lemma}

\theoremstyle{definition}
\newtheorem{definition}[theorem]{Definition}
\newtheorem*{example}{Example}
\newtheorem*{fact}{Fact}

\theoremstyle{remark}
\newtheorem{case}{Case}

%% Custom math commands
\newcommand{\E}{\mathbb{E}}
\newcommand{\1}{\mathds{1}}
\newcommand{\p}{\mathbb{P}}
\renewcommand{\H}{\mathcal{H}}
\newcommand{\T}{\mathcal{T}}
\newcommand{\ind}{\perp\!\!\!\perp}

%% File-specific additions
\usepackage{enumitem}

\newcommand{\X}{X}
\newcommand{\A}{A}
\newcommand{\Sx}{S}
\newcommand{\tauf}{\tau}
\newcommand{\cf}{c}
\newcommand{\gf}{g}
\newcommand{\muf}{\mu}
\newcommand{\mzero}{m_{0}}
\newcommand{\Spot}{S}
\newcommand{\RCT}{\mathrm{RCT}}
\newcommand{\RWD}{\mathrm{RWD}}
\newcommand{\RMST}{\mathrm{RMST}}
\newcommand{\SD}{\mathrm{SD}}
\newcommand{\AF}{\mathrm{AF}}

\title{Causal survival estimands under the AFT data-fusion model\\
with HDP error}
\author{Tijn Jacobs}
\date{\today}

\setlength{\parindent}{0pt}
\setlength{\parskip}{0.8em}


\begin{document}
\maketitle

\section{Setup}

We work under the AFT decomposition
\begin{equation}
  \log T_{i} = \mzero(\X_{i}, \Sx_{i})
            + \A_{i}\,\tauf(\X_{i})
            + (1 - \Sx_{i})\,\A_{i}\,\cf(\X_{i})
            + \sigma\,\varepsilon_{i},
  \qquad \varepsilon_{i} \sim P,
  \label{eq:model}
\end{equation}
with $P$ assigned a hierarchical Dirichlet process (HDP) prior with Gaussian
kernels and atoms shared across the two data sources. Conditional on the
sampler state, the implied error CDF in source $s$ is
\begin{equation}
  \Phi_{s}(z) = \sum_{k=1}^{K} \pi_{s, k}\,
                \Phi^{\mathrm{N}}\!\left(\frac{z - \mu_{k}}{\tau_{k}}\right),
  \label{eq:hdp-cdf}
\end{equation}
where $\Phi^{\mathrm{N}}$ is the standard normal CDF,
$\{(\mu_{k}, \tau_{k})\}_{k=1}^{K}$ are the cluster atoms (shared across
sources via the top-level DP), and $\pi_{s,k}$ are the source-specific
weights.

The treatment-specific AFT shift at target population $s_{t}$ is
\begin{equation}
  \eta_{a}(x, s_{t}) := \mzero(x, s_{t})
                    + a\,\tauf(x)
                    + (1 - s_{t})\,a\,\cf(x),
  \qquad a \in \{0, 1\}.
\end{equation}
The potential survival function is
\begin{equation}
  \Spot_{a}(t \mid x, s_{t})
  = 1 - \Phi_{s_{t}}\!\left(\frac{\log t - \eta_{a}(x, s_{t})}{\sigma}\right),
  \label{eq:potential-S}
\end{equation}
and we collect causal estimands as functionals of
$\{\Spot_{0}(\cdot \mid x, s_{t}), \Spot_{1}(\cdot \mid x, s_{t})\}$.

For the closed-form expressions below it is convenient to define, for each
component $k$ of the HDP mixture,
\begin{equation}
  \xi_{a, k}(x, s_{t}) := \eta_{a}(x, s_{t}) + \sigma\mu_{k},
  \qquad
  \omega_{k} := \sigma\tau_{k},
\end{equation}
so that the per-component survival function under the Gaussian-kernel HDP
is lognormal:
$1 - \Phi^{\mathrm{N}}\bigl((\log t - \xi_{a, k})/\omega_{k}\bigr)$.

\section{The survival difference $\Delta_{\SD}$}

\subsection{Definition}

The survival (risk) difference at time $t$ is
\begin{equation}
  \Delta_{\SD}(t; x, s_{t})
  := \Spot_{1}(t \mid x, s_{t}) - \Spot_{0}(t \mid x, s_{t}).
\end{equation}

\subsection{Closed form under the HDP}

\begin{equation}
  \boxed{\;
    \Delta_{\SD}(t; x, s_{t})
    = \sum_{k=1}^{K} \pi_{s_{t}, k}\!\left[
        \Phi^{\mathrm{N}}\!\left(\frac{\log t - \xi_{0, k}(x, s_{t})}{\omega_{k}}\right)
      - \Phi^{\mathrm{N}}\!\left(\frac{\log t - \xi_{1, k}(x, s_{t})}{\omega_{k}}\right)
    \right].
  \;}
  \label{eq:SD-closed}
\end{equation}

\subsection{Interpretation}

$\Delta_{\SD}(t; x, s_{t})$ is the absolute increase in the probability of
surviving past time $t$ attributable to assigning treatment versus control
for individuals with covariates $x$ in target population $s_{t}$. It lives
on the probability scale, is bounded between $-1$ and $1$, and has a direct
clinical reading: ``treatment increases the chance of surviving to year $t$
by $\Delta_{\SD}$ percentage points.''

This is the absolute risk reduction (ARR), or its survival-scale analogue.
It is the estimand most directly tied to clinical decision-making at a fixed
horizon and is the natural quantity for computing the number-needed-to-treat
$1/|\Delta_{\SD}(t; x, s_{t})|$ when the contrast is favourable
\citep{Cook1995}.

\paragraph{When to use.} When a single clinically meaningful horizon
exists --- for example, ``5-year survival in early-stage cancer'' or
``1-year survival post-transplant'' --- the SD is the most interpretable
summary. It depends on the error distribution $\Phi_{s_{t}}$ through
\eqref{eq:SD-closed} and is therefore sensitive to misspecification of the
tail when $t$ extends beyond the RCT follow-up; in the present framework
that sensitivity is absorbed by the HDP error.

\paragraph{Limitations.} $\Delta_{\SD}(t; x, s_{t})$ is a pointwise contrast
at a single $t$ and can give a misleading picture when survival curves
cross or when the gain is concentrated outside the chosen window. It does
not summarise the full survival benefit. For these reasons, regulatory and
HTA guidance recommends pairing the SD with an integrated summary such as
the RMST \citep{Royston2013}.

\section{The restricted mean survival time difference $\Delta_{\RMST}$}

\subsection{Definition}

The RMST under treatment $a$ at horizon $t^{*}$ is
\begin{equation}
  \mu_{a}(t^{*}; x, s_{t}) := \int_{0}^{t^{*}}\!
  \Spot_{a}(u \mid x, s_{t})\,du,
\end{equation}
and the RMST contrast is
\begin{equation}
  \Delta_{\RMST}(t^{*}; x, s_{t}) := \mu_{1}(t^{*}; x, s_{t}) - \mu_{0}(t^{*}; x, s_{t}).
\end{equation}

\subsection{Closed form under the HDP}

The per-component lognormal RMST has the standard expression
\begin{equation}
  \mu_{a}^{(k)}(t^{*}; x, s_{t})
  = e^{\,\xi_{a, k} + \omega_{k}^{2}/2}\,
    \Phi^{\mathrm{N}}\!\left(\frac{\log t^{*} - \xi_{a, k} - \omega_{k}^{2}}{\omega_{k}}\right)
  + t^{*}\!\left[1 - \Phi^{\mathrm{N}}\!\left(\frac{\log t^{*} - \xi_{a, k}}{\omega_{k}}\right)\right],
\end{equation}
and the contrast is the source-specific mixture
\begin{equation}
  \boxed{\;
    \Delta_{\RMST}(t^{*}; x, s_{t})
    = \sum_{k=1}^{K} \pi_{s_{t}, k}\bigl[
        \mu_{1}^{(k)}(t^{*}; x, s_{t}) - \mu_{0}^{(k)}(t^{*}; x, s_{t})
      \bigr].
  \;}
  \label{eq:RMST-closed}
\end{equation}

\subsection{Interpretation}

$\Delta_{\RMST}(t^{*}; x, s_{t})$ is the expected gain in survival time
within the window $[0, t^{*}]$ attributable to treatment, for individuals
with covariates $x$ in target population $s_{t}$. It lives on the
\emph{time scale}, in the same units as the outcome, and has a direct
reading: ``treatment yields $\Delta_{\RMST}$ additional units of expected
life within the first $t^{*}$ units.''

The RMST has been advanced as a principal alternative to the hazard ratio,
particularly when proportional hazards is implausible
\citep{Royston2011, Royston2013, Uno2014}. Its appeal is threefold: it has
a transparent units-of-time interpretation; it is well-defined under
non-proportional hazards including crossing survival curves; and it
integrates the full survival benefit within the window rather than
reporting a contrast at a single point. \citet{Uno2014} review its use in
oncology trials and provide guidance on the choice of $t^{*}$.

\paragraph{When to use.} The RMST is the natural headline estimand under
the long-term-extrapolation goal of this work. Reporting
$\Delta_{\RMST}(t^{*}; x, s_{t})$ at multiple horizons --- e.g.\ at
$t^{*}$ inside the trial window, at $t^{*}$ just past it, and at
$t^{*}$ far beyond it --- traces how the survival benefit evolves with
horizon and exposes the contribution of the HDP-tail to the inference.

\paragraph{Choice of $t^{*}$.} The horizon $t^{*}$ should be chosen
ex ante and motivated by clinical relevance or by regulatory context.
\citet{Royston2013} recommend $t^{*}$ at or near the largest follow-up
time at which a non-negligible proportion of the at-risk set remains
under observation. In the present setting that recommendation is too
conservative: the central methodological point is to extrapolate beyond
the RCT follow-up, so $t^{*}$ should deliberately exceed it.

\paragraph{Limitations.} The RMST depends on $t^{*}$, and substantive
conclusions can shift when $t^{*}$ is varied. Beyond the RWD follow-up,
$\Delta_{\RMST}$ inherits the full extrapolation uncertainty of the
posterior and should be reported with appropriate credible intervals.

\section{The acceleration factor and the CATE (log AF)}
\label{sec:AF}

Going forward we focus on the acceleration factor $\AF$ and the
conditional log acceleration factor $\tauf$. These are the same object on
two scales --- $\tauf(x) = \log \AF(x)$ --- and they are the natural
causal estimands under the AFT structural model. They are interpretable,
horizon-free, and (as we will see) free of the built-in selection bias
that plagues the hazard ratio.

\subsection{Definition}

Under the AFT model \eqref{eq:model}, the treatment-versus-control contrast
at fixed covariates $x$ is a constant additive shift on the log-time scale.
Equivalently, on the time scale it is a multiplicative \emph{acceleration
factor}. The causal acceleration factor at $x$ is
\begin{equation}
  \AF_{\mathrm{causal}}(x) := \exp\bigl\{\tauf(x)\bigr\}, \qquad
  \tauf(x) = \log \AF_{\mathrm{causal}}(x),
  \label{eq:AF-causal}
\end{equation}
i.e.\ the conditional log acceleration factor is the CATE on the log-time
scale.

The implied relationship between potential survival functions is the
quantile-shifting identity
\begin{equation}
  \Spot_{1}(t \mid x) = \Spot_{0}\bigl(\AF_{\mathrm{causal}}(x)\cdot t \mid x\bigr),
  \label{eq:AF-shift}
\end{equation}
so that survival under treatment at time $t$ equals survival under control
at the rescaled time $\AF_{\mathrm{causal}}(x)\,t$ \citep{CoxOakes1984,
Wei1992, KalbfleischPrentice2002}. The acceleration factor parameterises the entire
contrast between the two survival curves through a single multiplicative
time-scale parameter.

\paragraph{Causal versus observed contrast in the RWD.}
The corresponding \emph{observed} contrast in the RWD is
$\AF_{\mathrm{obs, RWD}}(x) = \exp\{\tauf(x) + \cf(x)\}$. The discrepancy
$\AF_{\mathrm{obs, RWD}}(x) / \AF_{\mathrm{causal}}(x) = \exp\{\cf(x)\}$
quantifies the multiplicative bias on the time scale due to unmeasured
confounding in the RWD. By transportability of potential outcomes
(Assumption~4 of the framework document), $\AF_{\mathrm{causal}}(x)$ is the
same in the RCT and RWD populations: it is a property of the causal
mechanism, not of the source.

\subsection{Closed form and HDP-independence}

The acceleration factor is a direct functional of the BART forest for
$\tauf$ and does not involve the error distribution:
\begin{equation}
  \boxed{\;\AF_{\mathrm{causal}}(x) = \exp\{\tauf(x)\}.\;}
\end{equation}
At each MCMC iteration, $\exp\{\tauf^{(r)}(x)\}$ yields a posterior draw of
$\AF_{\mathrm{causal}}(x)$ without reference to the cluster atoms or
weights of the HDP. Equivalently, posterior inference for $\tauf(x)$
\emph{is} posterior inference for the log acceleration factor. In the
context of long-term extrapolation (goal G2 of the project), this means
the acceleration factor is the one causal estimand that is unaffected by
the HDP tail and the extrapolation uncertainty it carries, in contrast to
$\Delta_{\SD}$ and $\Delta_{\RMST}$.

\subsection{Causal validity in the presence of unmeasured heterogeneity}
\label{subsec:af-causal-validity}

The clinical and methodological appeal of the acceleration factor rests on
its causal interpretation, which has recently been formalised by
\citet{Brathovde2024} under a structural-causal-model formulation of the
AFT. We summarise their result here because it is central to the
positioning of this work.

\citet{Brathovde2024} consider structural assignments for the potential
outcomes
\begin{equation}
  T^{0} := \exp\{f_{0}(U_{0}) + N_{T}\}, \qquad
  T^{a} := T^{0}/\exp\{f_{1}(U_{1}, a)\},
  \label{eq:scm}
\end{equation}
where $U_{0}$ represents \emph{frailty} (unmeasured heterogeneity in the
baseline hazard), $U_{1}$ represents \emph{effect heterogeneity}
(unmeasured modification of the treatment effect), and $N_{T}$ is residual
noise with $N_{T} \perp\!\!\!\perp (U_{0}, U_{1})$. The conditional causal
acceleration factor at $(U_{0}, U_{1}) = (u_{0}, u_{1})$ is then
$\theta_{c}(u_{0}, u_{1}) = \exp\{f_{1}(u_{1}, a)\}$ and is invariant in
$u_{0}$: frailty does not affect the causal AF interpretation
\citep{Keiding1997, Aalen2015}.

Their key identification result \citep[Theorem~3.1]{Brathovde2024} states
that, under no confounding, the \emph{observed} acceleration factor
\begin{equation}
  \theta_{m}(t) := \frac{1}{t}\,
  S^{-1}_{T \mid A=0}\!\bigl(S_{T \mid A=a}(t)\bigr)
\end{equation}
equals the marginal causal acceleration factor at every $t$ in the support
of the survival times. Equivalently, the AFT estimand identifies the
causal contrast even when both frailty and effect heterogeneity are
present in the population --- a property the hazard ratio does \emph{not}
share \citep[see also][]{hernan2010, MartinussenVansteelandtAndersen2020,
Post2024b}.

Under unmeasured confounding, $\theta_{m}$ no longer identifies $\theta$
directly. The standard remedy --- conditioning on a sufficient adjustment
set $L$ --- yields a conditional observed acceleration factor that equals
the conditional causal AF, and integrating over $F_{L}$ recovers $\theta$
\citep{Brathovde2024}. In our setup the same logic applies to $X$ in the
RCT (Assumption~2); in the RWD, the residual bias is encoded in $\cf(x)$,
and $\AF_{\mathrm{causal}}(x) = \exp\{\tauf(x)\}$ is identified from the
joint RCT$+$RWD data by the decomposition \eqref{eq:model}.

\subsection{The acceleration factor versus the hazard ratio}

The hazard ratio is the conventional summary in survival analysis, but it
is now well-understood to be a problematic causal estimand under
unmeasured heterogeneity. The literature on this point is mature:

\begin{itemize}[leftmargin=*]
  \item \citet{hernan2010} (``The hazards of hazard ratios'') argues that
        the HR is selection-biased because it conditions on survival,
        so that the survivor populations being compared at time $t > 0$
        no longer have exchangeable distributions of unmeasured
        prognostic variables.
  \item \citet{Aalen2015} formalise this: even in a randomised
        trial, the HR estimated by Cox regression does not in general
        coincide with the causal hazard ratio when frailty is present.
  \item \citet{MartinussenVansteelandtAndersen2020} review the subtleties
        and give worked examples in which the HR moves over time as a
        direct consequence of selection on frailty.
  \item \citet{Post2024b} prove that the Cox estimand differs from the
        causal effect of interest under frailty and effect heterogeneity,
        and that the AF maintains the causal interpretation in both
        cases.
\end{itemize}

The acceleration factor avoids this problem because it is defined on the
survival scale (a contrast between full marginal distributions) rather
than on the hazard scale (a contrast between conditional rates among
survivors). It does not require conditioning on survival, and is therefore
not subject to built-in selection bias \citep{Brathovde2024}. For a
project whose central methodological aim is to combine an RCT with an
real-world data --- both of which carry unmeasured heterogeneity ---
the acceleration factor is the most defensible causal target.

\subsection{Quantile equivariance}

Under the AFT model, the acceleration factor equals the ratio of
survival quantiles at every level:
\begin{equation}
  \frac{\mathrm{Quantile}_{q}\bigl(T(1) \mid x\bigr)}
       {\mathrm{Quantile}_{q}\bigl(T(0) \mid x\bigr)}
  = \AF_{\mathrm{causal}}(x)
  \qquad \text{for all } q \in (0, 1),
\end{equation}
a property called \emph{quantile equivariance} \citep{Wei1992,
Hougaard1999}. Three operational consequences:
\begin{enumerate}[leftmargin=*]
  \item \emph{Median survival ratio.} The most clinically natural
        instance: the AF is the ratio of treated to control median
        survival times. ``Treatment increases median survival by
        $(\AF_{\mathrm{causal}}(x) - 1)\times 100\%$'' is a direct reading.
  \item \emph{Any quantile, same number.} The 25th-percentile ratio, the
        90th-percentile ratio, and the median ratio all equal
        $\AF_{\mathrm{causal}}(x)$. The summary is therefore robust to the
        choice of survival quantile, in contrast to time-point survival
        differences.
  \item \emph{Direct comparability of curves.} Equation
        \eqref{eq:AF-shift} states that the treated and control survival
        curves differ by a horizontal stretch on the log-time axis, with
        stretch factor $\AF_{\mathrm{causal}}(x)$. This makes
        visualisation of the effect straightforward: overlay the two
        curves on a log-time axis, and the AF is the horizontal
        displacement.
\end{enumerate}

\subsection{Conditional, marginal, and time-varying AFs under heterogeneity}

In the present framework $\tauf(x) = \log\AF_{\mathrm{causal}}(x)$ is a
\emph{conditional} log acceleration factor, conditional on the observed
covariates $X = x$. Two subtleties about marginalising over the
population deserve attention.

\paragraph{Time-invariance is a property of the conditional AF.}
Under \eqref{eq:model}, $\AF_{\mathrm{causal}}(x)$ does not depend on $t$:
treatment is a single time-scale acceleration at each $x$. This is the
AFT shift assumption, and it is what makes extrapolation beyond the RCT
follow-up tractable.

\paragraph{Marginal AFs can be time-varying even when conditional AFs
are not.}
\citet{Brathovde2024} show, however, that even when each individual has
a time-invariant causal acceleration factor, the \emph{marginal} (i.e.\
population-level) AF can be time-varying when there is heterogeneity in
the AF across the population. This is a consequence of the marginalisation
itself: at early survival times, the population contains a different mix
of low-AF and high-AF individuals than at late survival times, because
survival probabilities at any fixed $t$ differ across the AF strata. The
practical consequence is that a constant individual AF cannot be
distinguished from a time-varying homogeneous AF on the basis of
population-level data alone. Conditioning on $X$ (and ideally on as much
of the relevant heterogeneity as can be measured) is therefore central.

For our reporting strategy, this implies:
\begin{itemize}[leftmargin=*]
  \item The headline summaries should be on the \emph{conditional} AF
        $\AF_{\mathrm{causal}}(x)$, e.g.\ via posterior projections of
        $\tauf(x)$ onto a small basis (cf.\ Woody--Carvalho--Murray
        projection summaries in the framing document).
  \item Population-averaged summaries should distinguish
        $\mathbb{E}[\AF_{\mathrm{causal}}(X)]$ from
        $\exp\{\mathbb{E}[\tauf(X)]\}$; the two coincide only under
        degenerate heterogeneity.
  \item Time-varying acceleration factor models
        \citep{Crowther2023, pang2021flexible} are available if the AFT shift
        assumption is itself questionable; we view these as a sensitivity
        analysis rather than a default.
\end{itemize}

\subsection{Clinical use and reporting}

The acceleration factor has a direct clinical reading that requires no
training in survival analysis. ``$\AF_{\mathrm{causal}}(x) = 1.5$ means
that, for a patient with profile $x$, treatment extends time-to-event by
50\% relative to control'' is communicable, transparent, and on the same
scale as the outcome itself. We recommend the following reporting
practices:

\begin{enumerate}[leftmargin=*]
  \item Report $\AF_{\mathrm{causal}}(x)$ on the multiplicative scale
        (e.g., $\AF = 1.5$) and the CATE $\tauf(x)$ on the log scale
        side by side. The log scale is preferred for inference (additive,
        symmetric posterior); the multiplicative scale is preferred for
        clinical communication.
  \item Report subgroup or covariate-conditional AFs via posterior
        projection of $\tauf(x)$ onto an interpretable basis, with
        credible intervals derived from the projection posterior.
  \item Pair AFs with at least one absolute summary --- the SD at a
        clinically meaningful $t$, or $\Delta_{\RMST}(t^{*})$ at a
        clinically meaningful $t^{*}$ --- to communicate the
        \emph{magnitude} of the absolute benefit. A 50\% time-scale
        acceleration on a baseline median of one month and on a baseline
        median of ten years have very different clinical implications.
  \item Where extrapolation is concerned (goal G2), prefer the AF as the
        headline summary at long horizons: it is the estimand least
        contaminated by tail uncertainty.
\end{enumerate}

\paragraph{Caveat: AFT shift as a structural assumption.}
The acceleration-factor framework depends on the AFT shift structure of
\eqref{eq:model}: that the treatment effect at fixed $x$ is a single
multiplicative acceleration of the time scale that does not vary with
$t$. When this is doubtful (waning treatment effects, mode-of-failure
shifts, late-emerging benefits), one can either (i) accept the AFT
summary as a parsimonious approximation and discuss in sensitivity, or
(ii) generalise to a time-varying acceleration factor
$\AF(x, t)$ via the flexible parametric approach of
\citet{Crowther2023} or the time-dependent AFT extensions of
\citet{pang2021flexible}. The latter generalisation comes at the cost of giving
up the single-number summary that motivates the AF in the first place.

\paragraph{Note on time-varying treatment regimes.}
For studies with time-varying treatments and treatment--confounder
feedback, the structural-causal-AFT framework of
\citet{HernanColeStrAFT2005} (using $g$-estimation) is required. Our
setting --- a single point-in-time treatment assignment --- falls within
the scope of \citet{Brathovde2024}'s simpler structural AFT and does not
require this generalisation.

\section{Population-averaged versions}

Each of the three estimands has a population-averaged counterpart obtained
by marginalising over the covariate distribution of the target population:
\begin{align}
  \Delta_{\SD}(t; s_{t})
  &= \mathbb{E}\bigl[\Delta_{\SD}(t; \X, s_{t}) \mid \Sx = s_{t}\bigr], \\
  \Delta_{\RMST}(t^{*}; s_{t})
  &= \mathbb{E}\bigl[\Delta_{\RMST}(t^{*}; \X, s_{t}) \mid \Sx = s_{t}\bigr], \\
  \AF_{\mathrm{causal}}
  &= \mathbb{E}\bigl[\AF_{\mathrm{causal}}(\X)\bigr]
   \quad \text{or} \quad
   \exp\bigl(\mathbb{E}[\tauf(\X)]\bigr),
\end{align}
where the last definition distinguishes the population-averaged AF (which
preserves the multiplicative interpretation pointwise but is no longer a
ratio of population means) from the AF of the population-averaged log-time
shift (which is). The two coincide only under degenerate $\tauf(\X)$.
Pointwise reporting via the posterior projection summaries is generally
preferable to a single marginal number unless heterogeneity is small.

The Monte Carlo estimator is in each case the empirical average over
evaluation rows of the corresponding posterior draw, with Bayesian-bootstrap
weights if covariate-distribution uncertainty is to be propagated.

\section{Summary}

\begin{itemize}[leftmargin=*]
  \item $\Delta_{\SD}(t; x, s_{t})$ --- absolute survival-probability
        gain at fixed $t$. Probability scale; depends on HDP tail.
        Natural for fixed clinical horizons and number-needed-to-treat.
  \item $\Delta_{\RMST}(t^{*}; x, s_{t})$ --- expected life-time gain
        within $[0, t^{*}]$. Time scale; depends on HDP tail.
        Natural headline estimand under the long-term extrapolation goal.
  \item $\AF_{\mathrm{causal}}(x) = \exp\{\tauf(x)\}$ --- multiplicative
        time-scale acceleration. Horizon-free; HDP-free; assumes
        the AFT shift structure. Natural complement summarising treatment
        heterogeneity in a single covariate-specific number.
\end{itemize}

\bibliographystyle{apalike}
\bibliography{../general/references}

\end{document}
